% !PRINTED none

\begin{center}
\gesperrt{Eighth Section.}

\textbf{\gesperrt{The Polyhedral Groups.}}

\rule{3cm}{0.4pt}

\S.\ 55.
\end{center}

\begin{center}
\gesperrt{The cyclic groups and the dihedral groups.}
\end{center}

We now proceed to apply the general principles to the actual formation of the
various polyhedral groups, and thus first to show that the kinds of these groups
recognised as possible in \S.\,52 all exist.

For the cyclic groups of degree $n$ we have only two
systems of conjugate poles, each of which contains only one $(n-1)$-fold
pole. These two poles must therefore be the common
poles of all substitutions of the group, and by
\S.\,51, 2. may be taken as $0$ and $\infty$. We then have
only the two linear fundamental forms
\weq{1}{f_1 = x_1, \quad f_2 = x_2.}
All substitutions of the group are multiplicative, and the
multiplier must be an $n$th root of unity. They must therefore
have the form:
\weq{2}{x, \varepsilon x, \varepsilon^2 x, \ldots, \varepsilon^{n-1} x,}
if $\varepsilon$ is a primitive $n$th root of unity.

We thus obtain in fact for every $n$ a cyclic
group, which we shall denote by $C_n$:
\weq{3}{C_n = \begin{pmatrix} \varepsilon^r, & 0 \\ 0, & 1 \end{pmatrix}, \quad r = 0, 1, \ldots, n-1,}
or written with the determinant $1$:
\weq{4}{C_n = \begin{pmatrix} \varepsilon^{1/2\,r}, & 0 \\ 0, & \varepsilon^{-1/2\,r} \end{pmatrix}.}
